\clearpage
\section{Normative XOP-R1 grammar and operator registry}
\label{sec:exact-registry}
This chapter embeds the literal, type, wire and operator contract from
\path{formal/OPERATOR_CATALOG.md}. It supplies the actual fields and numbered
operators needed to read the proposed profile without another file. The
mathematical denotations and conditional proofs are in
Section~\ref{sec:exact-core}; the attitude application is in
Section~\ref{sec:exact-madgwick}. This is a specification, not an implemented
parser, exact evaluator or proof checker. Unassigned tags and opcodes are
rejected; the new grammar does not reinterpret the preserved R10 format.

\begingroup
\newcommand{\xo}[1]{\texttt{\detokenize{#1}}}
\setlength{\tabcolsep}{4pt}
\renewcommand{\arraystretch}{1.11}

\subsection{Sorts, outcomes and canonical byte primitives}
The scalar embeddings are explicit and value preserving:
$\xo{Int}\hookrightarrow\xo{Rat}\hookrightarrow\xo{Alg}
\hookrightarrow\xo{RealTerm}$. Booleans and finite words are separate types.
\xo{RealTerm} carries finite syntax with a possibly partial real denotation;
it is not a literal for an unspecified arbitrary real number. Structural types
carry their shapes. A \xo{Dim} scalar additionally carries seven rational SI
exponents. A \xo{Frame} annotation encodes nominal coordinate identity; prose
alone cannot supply it. Angle conventions remain declared interface metadata.

\begin{longtable}{@{}P{.24\linewidth}P{.71\linewidth}@{}}
\toprule Outcome & Meaning\\\midrule\endfirsthead
\toprule Outcome & Meaning (continued)\\\midrule\endhead
\bottomrule\endfoot
\xo{VALUE} & An exact literal with evidence of its construction.\\
\xo{TERM} & Retained typed syntax, assumptions and outstanding obligations.\\
\xo{UNDEFINED} & A violated domain condition has been established.\\
\xo{UNKNOWN} & A required comparison, domain or existence claim is undecided.\\
\xo{MISSING_INPUT} & A required dependency or sample is unavailable.\\
\xo{RESOURCE_LIMIT} & The declared memory, time or proof-search limit is reached.\\
\end{longtable}
These outcomes are separate from scalar values. Unknown or exhausted resources
mean neither false nor zero. A retained term is not a proved real value.
\xo{COMPARE} is a decision-service interface returning justified
\xo{LT}, \xo{EQ}, \xo{GT}, or \xo{UNKNOWN}; it is not an arithmetic oracle.

Concatenation in the grammar has no implicit alignment. The byte primitives are:
\begin{longtable}{@{}P{.22\linewidth}P{.73\linewidth}@{}}
\toprule Primitive & Encoding\\\midrule\endfirsthead
\toprule Primitive & Encoding (continued)\\\midrule\endhead
\bottomrule\endfoot
\xo{U(n)} & Minimal unsigned base-128 integer: seven payload bits per octet,
least significant digit first; high bit one on every nonterminal octet and zero
on the last. Zero is one zero octet. A multi-octet encoding cannot end in a zero
payload. Counts have no format-level 64-bit ceiling.\\
\xo{O(b)} & One octet, $0\le b\le255$.\\
\xo{B(bytes)} & \xo{U(length)} followed by exactly those octets.\\
\xo{Seq_T(items)} & \xo{U(count)} followed by the items in order, each using
schema $T$.\\
\xo{Z(z)} & \xo{O(sign) U(limb_count)}, then that many eight-octet,
little-endian limbs; canonical sign and leading-limb rules below.\\
\xo{Q(q)} & \xo{Z(numerator) Z(denominator)}, with positive denominator and
reduced fraction.\\
\xo{Name} & A \xo{B} byte string containing only printable ASCII octets
\xo{0x21..0x7e}; nonempty and case sensitive. Free-form provenance is opaque
\xo{B} data and has no name-binding semantics.\\
\xo{ProfileId} & \xo{Name(namespace) U(major_version) U(minor_version) Name(label)}.
Complete fields form the nominal type/binding identity; they are not executable prose.\\
\end{longtable}
An implementation may announce finite count, width, depth and memory limits.
A larger valid mathematical record is then unsupported, not truncated or partly
accepted. Check lengths before allocation and retain exact count arithmetic.

\subsection{Literal payloads and type tags}
For integer limbs $a_i$ in base $\beta=2^{64}$,
$z=(-1)^s\sum_{i=0}^{k-1}a_i\beta^i$, with $0\le a_i<\beta$.
Zero is $s=0,k=0$; otherwise $a_{k-1}\ne0$. A rational $p/q$ requires
$q>0$ and $\gcd(|p|,q)=1$; rational zero is $0/1$.

\begin{longtable}{@{}P{.23\linewidth}P{.72\linewidth}@{}}
\toprule Literal type & Exact payload and admission rules\\\midrule\endfirsthead
\toprule Literal type & Exact payload and admission rules (continued)\\\midrule\endhead
\bottomrule\endfoot
\xo{Bool} & \xo{O(value)}, with value zero or one.\\
\xo{Int} & \xo{Z(value)}, with the canonical integer form above.\\
\xo{Rat} & \xo{Q(value)}, with the canonical rational form above.\\
\xo{Alg} & \xo{Seq_Z(coefficients) U(root_index)}. Coefficients run from
constant to highest degree. The polynomial has degree at least one, is primitive
and irreducible over $\mathbb Q$, and has positive leading coefficient. The
zero-based index selects a distinct real root in increasing order.\\
\xo{Bits(w)} & Exactly $\lceil w/8\rceil$ octets; bit zero is the low bit of
the first octet. Unused high bits in the final octet are zero. Width is in the
type, not repeated in this payload.\\
\xo{Dim(T,d)} & Exactly the scalar payload of $T$; the seven exponents already
occur in its type. No extra numeric field or raw \xo{RealTerm} literal is
authorized by the wrapper.\\
Tuple, vector, matrix & Ordered \xo{Seq_Literal} components with type/shape
agreement; each component retains its own type.\\
\xo{Frame} & The underlying Vec/Mat literal payload, with the complete frame
annotation retained in \xo{Type}; encoding and sharing cannot erase it.\\
\end{longtable}
An algebraic isolating interval is a witness, not an alternative value spelling.
Exact factorization, minimal-polynomial selection, root isolation and selected
root identification remain implementation obligations. A resultant alone does
not select the intended root. Raw finite IEEE bits can be converted to their
exact dyadic value; this does not recover a pre-rounding physical quantity.
Signed-zero payloads remain distinct as bits but both denote rational zero.
NaN and infinity have no real conversion.

Each type is \xo{U(tag)} followed by the fields in this table:
\begin{longtable}{@{}P{.08\linewidth}P{.22\linewidth}P{.65\linewidth}@{}}
\toprule Tag & Type & Remaining fields\\\midrule\endfirsthead
\toprule Tag & Type & Remaining fields (continued)\\\midrule\endhead
\bottomrule\endfoot
0 & \xo{Bool} & None.\\
1 & \xo{Int} & None.\\
2 & \xo{Rat} & None.\\
3 & \xo{Alg} & None.\\
4 & \xo{RealTerm} & None.\\
5 & \xo{Bits} & \xo{U(w)}, $w>0$.\\
6 & \xo{Tuple} & \xo{Seq_Type(component_types)}.\\
7 & \xo{Vec} & \xo{U(n) Type(T)}, $n>0$.\\
8 & \xo{Mat} & \xo{U(rows) U(cols) Type(T)}; both sizes positive.\\
9 & \xo{Fn} & \xo{Seq_Type(parameter_types) Type(result)}.\\
10 & \xo{Dim} & \xo{Type(scalar)}, followed by exactly seven \xo{Q} exponents.\\
11 & \xo{Stream} & \xo{Type(element)}.\\
12 & \xo{Frame} & \xo{Type(base) ProfileId(source)}, plus \xo{ProfileId(target)}
when base is Mat. Base is an unannotated Vec or Mat.\\
\end{longtable}
\xo{Dim} accepts only an unwrapped numeric scalar; nested dimension wrappers
are rejected. A function's domain obligations remain part of its interface/body.
Function, stream and ungrounded real-term literals are rejected.

A framed vector's identity denotes its coordinate basis; a framed matrix's
source/target denote the input/output bases of its linear map. Matching nominal
identity does not prove physical alignment, origin or calibration. Axis/origin
realization and uncertainty are explicit template/dependency inputs. Qualified
namespaces must agree across the dependency closure.

Bare vectors/matrices are frame-free objects. Once an interface declares a
frame, its annotation is mandatory. \xo{COERCE} cannot add, change or remove it.
Literal input may declare a frame; a computed frame change requires an explicitly
typed \xo{TEMPLATE}, its full body and coordinate/domain obligations. Extracting
coordinates and relabelling the reassembled result is not an implicit conversion.
This also governs removal of an annotation for use in geometry.

Frame-aware addition/subtraction and dot/cross require identical vector frame
identities. Norm returns a scalar. A matrix $A\to B$ maps a vector in $A$ to $B$
through an explicit contraction template; composing $A\to B$ then $B\to C$
produces $A\to C$ in the corresponding product order. Transpose and inverse
swap source/target annotations under their algebraic domains; interpreting
transpose as inverse also requires orthonormality. Determinant refers to the
declared coordinate matrix, without a basis-independence assertion. Projection
retains its annotated source in node identity even when its result is scalar.
Sharing/memoization include complete annotated types and immutable bindings.

\subsection{Scopes, references and complete record order}
The following fields are concatenated in the displayed order:
\begin{lstlisting}
Scope = Seq_Node(nodes) Seq_Ref(exports)
Node = U(opcode) Type(result) Seq_Ref(operands) B(opcode_payload)
Literal(T) = Type(T) B(payload)

ASCII("XOPSEED1") U(1) U(0)
Seq_Dependency(dependencies)
Seq_Template(templates)
Seq_State(states)
Seq_Sample(samples)
Scope(main)
Seq_Export(exports)
Strategy(strategy)
B(provenance)
\end{lstlisting}
The magic is exactly eight ASCII octets; the two naturals specify major/minor
profile $(1,0)$. No trailing octets are allowed. Every opcode payload must decode
exactly under its selected schema, with no extra fields. Empty payload means
a zero-length \xo{B}; the stated result type must satisfy the operator rule.

A reference is \xo{U(tag)} followed by:
\begin{longtable}{@{}P{.08\linewidth}P{.24\linewidth}P{.63\linewidth}@{}}
\toprule Tag & Reference kind & Fields and scope rule\\\midrule\endfirsthead
\toprule Tag & Reference kind & Fields and scope rule (continued)\\\midrule\endhead
\bottomrule\endfoot
0 & Local node & \xo{U(index)}. A node operand refers only to an earlier node
in its own scope.\\
1 & Bound variable & \xo{U(de_bruijn_index)}. Zero is the nearest lexical binder.\\
2 & Old state & \xo{U(slot)}; main dynamic scope only.\\
3 & Current sample & \xo{U(slot)}; main dynamic scope only.\\
4 & Dependency export & \xo{U(dependency) U(export)}; complete bytes and export
interface must resolve.\\
5 & Current sample raw timestamp & \xo{U(slot)}; main dynamic scope only.
Returns the binding's exact \xo{Q(absolute_time)} as Rat in the slot's source
\xo{TimeConvention} coordinate.\\
6 & Current sample GPST timestamp & \xo{U(slot)}; main dynamic scope only.
Explicitly applies the declared time mapping and returns
$\mathrm{Dim}(\xo{Rat},(0,0,1,0,0,0,0))$ in GPST SI seconds.\\
\end{longtable}
A scope export or state root may use any already declared node in its complete
scope, but this does not authorize forward operands inside a node. Nested
scopes cannot reach an outer local node through an accidental numeric index:
captures enter as explicit binders. Binder order follows the payload; the last
parameter has de Bruijn index zero. Substitution is capture avoiding and shifts
indices beneath binders. Templates receive state/sample values as arguments.

Tags 3, 5 and 6 read the same immutable sample snapshot. Tag 5 exposes recorded
source time, without asserting that a custom coordinate is SI seconds. Tag 6
preserves the registered GPST-SI rational with its time units, or invokes the
complete custom mapping. Missing input and unmet/unresolved conversion domains
remain explicit non-value outcomes; no offset is guessed. The SI exponent order
is length, mass, time, electric current, thermodynamic temperature, amount of
substance, luminous intensity.

An SI-second integrator's time arguments are declared expressions, for example
\[
 t=\operatorname{SAMPLE\_GPST}(j)-t_{\rm epoch},\qquad
 h=\operatorname{SAMPLE\_GPST}(j)-t_{\rm previous},
\]
with matching time dimensions and GPST epoch conventions on every term.
Here \xo{SAMPLE_GPST} names reference tag 6, not an additional opcode.
Other integration coordinates require explicit typed mappings and domain/unit
hypotheses. Chronology and $h>0$ remain recurrence obligations. Initial roots and
pure template bodies exclude tags 5 and 6 as they exclude tag 3; a template
receives needed timestamps as arguments. Raw ticks, phase wraps and civil-clock
strings cannot silently supply these integration arguments.

\begin{lstlisting}
Dependency = Name(label) O(kind) B(canonical_seed_bytes)
Template = Name(label) Seq_Type(parameters) Type(result) Scope(body)
State = Name(label) Type(T) Ref(init) Ref(next)
Sample = Name(label) Type(T) TimeConvention
         O(mode) B(binding) B(provenance)
Row = Z(logical_step) Q(absolute_time) Literal(T)
Export = Name(label) U(main_export_index) O(observation)
TimeConvention = ProfileId(identity) O(kind) B(definition)
\end{lstlisting}
\textbf{Dependencies.} Kind zero imports a seed value export; kind one imports
a template-library seed. Content is a complete \xo{XOPSEED1} record, labels are
unique, and imports are finite and acyclic. Imported values must be closed and
stateless in their demanded graph. A snapshot may be closed literal/expression
data; a live stream must enter through an explicit sample. Importing a seed
neither executes its state machine nor silently selects a clock tick.

\textbf{Templates.} A body has exactly one export of its declared result type.
It contains no main-state/sample references and calls only earlier local or
explicitly imported templates. The encoded body fixes its meaning; a label alone
does not provide an equation or authorize recursion.

\textbf{State.} Initial and next roots refer to the main scope. Initial roots
are closed with respect to state and samples and must have proved values of
the declared type. Next roots read the same old-state/input snapshot. The state
table is the sole causal \xo{DELAY}; all successful next values commit together.
No instantaneous cycle or partial commit is defined. A required input/domain
failure remains an explicit outcome.

\textbf{Samples.} Mode zero binding is \xo{Seq_Row(rows)} with nonnegative,
strictly increasing logical-step indices. Mode one binding is
\xo{Name(external_port)}; every demanded step supplies a typed literal and
absolute timestamp under the named convention. Missing rows or ports yield
\xo{MISSING_INPUT}. No implicit interpolation, nearest-row choice, leap-second
conversion or future extrapolation is supplied. Repeated reads in one step bind
the same sample. The timestamp and sample identity are retained.

\textbf{Time convention.} Kind zero has empty definition and exactly the
registered identity \xo{("XOP-TIME",1,0,"GPST-SI")}. Its rational coordinate is
elapsed SI seconds on continuous GPS time from 1980-01-06 00:00:00 GPST; no UTC
conversion occurs. Kind one is custom, with definition
\xo{U(dependency) U(template) ProfileId(target)}. The dependency is a complete
template-library seed and the target is the registered GPST-SI identity.
The selected pure template maps unambiguous rational source time to rational
target time. Its mapping, interval, epoch and any leap/disambiguation tables
must be in its body or complete dependencies. Ambiguous civil readings need
additional binding. Repeated custom identities in the dependency closure must
resolve to identical canonical definition/template bytes. Ordering and elapsed
time claims require the mapping's monotonicity/domain proof, not version
matching alone.

\textbf{Exports.} Labels are unique and observation is zero or one. Stateful
equivalence preserves observation-one exports, sample identities and declared
output timestamps. Even unobserved exports remain part of the ordinary value
interface unless an explicit interface change is adopted.

Canonical encoding means one spelling for the same ordered scoped record, not
for every mathematically equivalent expression or sharing pattern. Dependency
hashes are lookup aids; complete canonical bytes establish required identity.
A digest-only open manifest is not this self-contained encoding.

\subsection{The one-bit plane envelope}
Let $C$ be a canonical semantic stream of positive length $L$. Pad it with zero
octets to $512\lceil L/512\rceil$ octets. Each block holds 64 little-endian
uint64 patterns $w_r$. Its plane words are
\[
 p_c=\sum_{r=0}^{63}\left((w_r\mathbin{\gg}c)\mathbin{\&}1\right)2^r,
 \qquad 0\le c<64.
\]
\begin{lstlisting}
ASCII("XOPPLAN1") U(L)
plane words: eight octets each, little endian,
             block-major then plane-major order
\end{lstlisting}
Exactly $512\lceil L/512\rceil$ payload octets are required. Apply the same
transpose to invert; reject nonzero decoded padding, take exactly $L$ octets,
then parse \xo{XOPSEED1}. Nonminimal lengths, redundant integer forms, invalid
references, nonzero padding and trailing bytes have no canonical decoding.
The transpose proof is in Section~\ref{sec:exact-core}. Compression is a separate
transport that must reconstruct this envelope exactly. A bit permutation does
not reduce its uncompressed information content or erase high limbs.

\subsection{Structural, numeric and Boolean opcodes}
In the tables, $S$ is a common underlying scalar sort reached by explicit
coercions. Dimension transformations below refine the scalar signatures;
matching $S$ does not require multiplication operands to have equal units.
$T$ is a complete type. Unless stated otherwise, payload is empty. Ordinary
numeric operands are strict in demanded subexpressions; \xo{IF} is lazy and a
lambda body is demanded only through application. A node table is a dependency
order, not a command to evaluate unused branches.

{\footnotesize
\begin{longtable}{@{}P{.23\linewidth}P{.24\linewidth}P{.48\linewidth}@{}}
\toprule ID / operator & Operands $\to$ result & Domain and payload\\\midrule\endfirsthead
\toprule ID / operator & Operands $\to$ result & Domain and payload (continued)\\\midrule\endhead
\bottomrule\endfoot
0 \xo{LITERAL} & None $\to$ literal type & Payload is that type's canonical literal payload.\\
1 \xo{COERCE} & Scalar $\to$ higher scalar & Only the stated injective scalar embeddings.\\
2 \xo{TUPLE} & Ordered values $\to$ tuple & Result component types match.\\
3 \xo{PROJECT} & Tuple/Vec/Mat $\to$ component & Payload \xo{U(index)}, in bounds. Matrix index is row-major $r\cdot\mathrm{cols}+c$ and returns one scalar.\\
4 \xo{VECTOR} & $n$ scalars $\to\mathrm{Vec}(n,S)$ & Positive size; common scalar sort and dimension.\\
5 \xo{MATRIX} & $rc$ scalars $\to\mathrm{Mat}(r,c,S)$ & Row-major order; result fixes shape.\\
6 \xo{LAMBDA} & Captures $\to$ Fn & Payload \xo{Seq_Type(parameters) Scope(body)}; one body export; capture binders precede parameter binders.\\
7 \xo{APPLY} & Function, arguments $\to$ result & Exact arity/types; substitute into the body and its domain.\\
8 \xo{TEMPLATE} & Arguments $\to$ template result & Payload \xo{O(location) U(index)}. Location 0 selects local template. Location 1 makes index a dependency and adds \xo{U(imported_template_index)}.\\
16 \xo{ADD} & $S,S\to S$ & Equal dimensions.\\
17 \xo{SUB} & $S,S\to S$ & Equal dimensions.\\
18 \xo{NEG} & $S\to S$ & Numeric scalar.\\
19 \xo{MUL} & $S,S\to S$ & Add dimension vectors.\\
20 \xo{DIV} & $S,S\to S$ & Denominator nonzero; Int result forbidden, requiring explicit Rat promotion; subtract dimensions.\\
21 \xo{IDIV} & Int, Int $\to$ Int & Nonzero divisor; Euclidean quotient.\\
22 \xo{MOD} & Int, Int $\to$ Int & Nonzero divisor $d$; $0\le r<|d|$.\\
23 \xo{POW_INT} & $S$, Int $\to S$ & Nonnegative exponent, or nonzero base and rational/higher result; $x^0=1$ by finite-product convention.\\
24 \xo{ABS} & $S\to S$ & Ordered real scalar.\\
25 \xo{SIGN} & $S\to$ Int & $-1,0,1$ according to exact sign; decision can remain unresolved.\\
26 \xo{MIN} & $S,S\to S$ & Matching dimensions; order evidence required for selection.\\
27 \xo{MAX} & $S,S\to S$ & Matching dimensions; order evidence required for selection.\\
28 \xo{FLOOR} & Dimensionless scalar $\to$ Int & Unique $n$ with $n\le x<n+1$; decision can remain unresolved.\\
29 \xo{CEIL} & Dimensionless scalar $\to$ Int & $-\lfloor-x\rfloor$.\\
30 \xo{ROOT} & Int $n$, scalar $x$ $\to$ Alg/RealTerm & $n\ge1$. Even $n$: $x\ge0$, nonnegative root. Odd $n$: unique real root. Divide dimensions by $n$.\\
31 \xo{SQRT} & Scalar $\to$ Alg/RealTerm & $x\ge0$, nonnegative root; explicit scalar promotion and halved dimensions.\\
32 \xo{EQ} & Compatible equal types $\to$ Bool & Exact value equality; no terminating general RealTerm decision is promised.\\
33 \xo{LT} & $S,S\to$ Bool & Ordered scalar; matching dimensions.\\
34 \xo{LE} & $S,S\to$ Bool & Ordered scalar; matching dimensions.\\
35 \xo{AND_BOOL} & Bool, Bool $\to$ Bool & Logical conjunction, not a word operation.\\
36 \xo{OR_BOOL} & Bool, Bool $\to$ Bool & Logical disjunction.\\
37 \xo{NOT_BOOL} & Bool $\to$ Bool & Logical negation.\\
38 \xo{IF} & Bool, $T,T\to T$ & Lazy: only the selected branch must be defined. An unknown predicate retains the conditional term.\\
40 \xo{FULL_ADDER} & Bool $a,b,c\to$ Bool pair & Sum bit and carry bit, equations below.\\
41 \xo{FULL_SUBTRACTOR} & Bool $a,b,k\to$ Bool pair & Difference bit and borrow bit, equations below.\\
48 \xo{BIT_AND} & Bits$(w)$ pair $\to$ Bits$(w)$ & Equal explicit widths.\\
49 \xo{BIT_OR} & Bits$(w)$ pair $\to$ Bits$(w)$ & Equal explicit widths.\\
50 \xo{BIT_XOR} & Bits$(w)$ pair $\to$ Bits$(w)$ & Equal explicit widths.\\
51 \xo{BIT_NOT} & Bits$(w)\to$ Bits$(w)$ & Complement inside the declared width only.\\
52 \xo{BIT_SHIFT} & Bits$(w)$, Int $k\to$ Bits$(w)$ & $k\ge0$; payload \xo{O(direction)}: 0 left, 1 logical right. Shifted-out bits are deliberately discarded.\\
53 \xo{POPCOUNT} & Bits$(w)\to$ Int & Sum of the $w$ Boolean bits.\\
54 \xo{BITS_TO_UINT} & Bits$(w)\to$ Int & $\sum_i b_i2^i$.\\
55 \xo{UINT_TO_BITS} & Int $\to$ Bits$(w)$ & $0\le x<2^w$; no implicit modular truncation.\\
56 \xo{INT_SHIFT} & Int, Int $k\to$ Int & $k\ge0$; payload direction 0 multiplies by $2^k$, 1 takes floor division by $2^k$.\\
57 \xo{IEEE64_TO_RAT} & Bits$(64)\to$ Rat & Finite IEEE binary64 pattern only; exact dyadic conversion below.\\
58 \xo{DECIMAL_TO_RAT} & Int mantissa, Int scale $\to$ Rat & Exact $\mathrm{mantissa}\cdot10^{-\mathrm{scale}}$, reduced.\\
\end{longtable}}
For IEEE conversion let $s$ be the sign bit, $e$ the unsigned 11-bit exponent
and $f$ the unsigned 52-bit fraction. The exact value is
\[
 \begin{cases}
 (-1)^s f\,2^{-1074},&e=0,\\
 (-1)^s(2^{52}+f)\,2^{e-1075},&1\le e\le2046.
 \end{cases}
\]
Exponent 2047 is outside the real conversion domain. Keep the raw operand if
signed-zero or other bit provenance must remain observable. Rational operations
reduce by gcd and normalize sign; algebraic operations preserve selected-root
identity. The source equation-language token \xo{1} means an all-ones word only
under its separately declared language adapter.

For Boolean carry $c$ and borrow $k$, the one-bit arithmetic operators are
\begin{align*}
 s&=a\oplus b\oplus c,&c'&=(a\land b)\lor(c\land(a\oplus b)),\\
 d&=a\oplus b\oplus k,&k'&=(\neg a\land b)\lor(k\land\neg(a\oplus b)),\\
 a+b+c&=s+2c',&a-b-k&=d-2k'.
\end{align*}
Complements here act on one Boolean bit. Chain carry/borrow toward more
significant bits, starting at zero. Multiply by placing $a_i\land b_j$ at
position $i+j$ and adding every shifted row with retained carries. To compare
padded magnitudes from high to low bits, start $(E,G)=(1,0)$ and update
simultaneously
$G'=G\lor(E\land a\land\neg b)$ and $E'=E\land\neg(a\oplus b)$.
Long division uses this exact comparison and subtraction. These are finite
definitions for every finite input; no executed bit test is asserted.

\subsection{Elementary, tensor, calculus and kernel opcodes}
Elementary functions have the real denotations and branch conventions in
Section~\ref{sec:exact-core}. They require explicit promotion to
\xo{RealTerm}; angle results are dimensionless with declared angle convention.

{\footnotesize
\begin{longtable}{@{}P{.23\linewidth}P{.24\linewidth}P{.48\linewidth}@{}}
\toprule ID / operator & Operands $\to$ result & Domain and payload\\\midrule\endfirsthead
\toprule ID / operator & Operands $\to$ result & Domain and payload (continued)\\\midrule\endhead
\bottomrule\endfoot
64 \xo{SIN} & RealTerm $\to$ RealTerm & Entire sine series; dimensionless argument.\\
65 \xo{COS} & RealTerm $\to$ RealTerm & Entire cosine series; dimensionless argument.\\
66 \xo{TAN} & $x\to\sin x/\cos x$ & $\cos x\ne0$.\\
67 \xo{EXP} & RealTerm $\to$ RealTerm & Entire exponential series.\\
68 \xo{LOG} & $x\to\int_1^x du/u$ & $x>0$.\\
69 \xo{ATAN} & $x\to\int_0^x du/(1+u^2)$ & Every real $x$.\\
70 \xo{ATAN2} & $y,x\to$ angle & $(x,y)\ne(0,0)$; range $(-\pi,\pi]$ with the declared quadrant cases.\\
71 \xo{ASIN} & $x\to$ angle & $-1\le x\le1$; inverse sine on $[-\pi/2,\pi/2]$.\\
72 \xo{ACOS} & $x\to$ angle & $-1\le x\le1$; inverse cosine on $[0,\pi]$.\\
73 \xo{PI} & None $\to$ RealTerm & $4\operatorname{atan}(1)$.\\
74 \xo{POW_REAL} & $x,y\to\exp(y\log x)$ & $x>0$; both arguments dimensionless. Distinct from integer powers and root branches.\\
80 \xo{DOT} & Vec$(n,S)$ pair $\to S$ & Finite sum of pairwise products.\\
81 \xo{CROSS} & Vec$(3,S)$ pair $\to$ Vec$(3,S)$ & Right-handed determinant in the declared frame.\\
82 \xo{NORM} & Vec$(n,S)\to$ Alg/RealTerm & Nonnegative square root of dot product; retains the vector dimension.\\
83 \xo{TRANSPOSE} & Mat$(r,c,S)\to$ Mat$(c,r,S)$ & Swap indices.\\
84 \xo{MATMUL} & Mat$(r,k,S)$, Mat$(k,c,S)\to$ Mat$(r,c,S)$ & Finite contraction over $k$.\\
85 \xo{DET} & Square matrix $\to S$ & Finite signed permutation sum.\\
86 \xo{INVERSE} & Mat$(n,n,S)\to$ Mat$(n,n,S)$ & Nonzero determinant; scalar sort closed under division.\\
87 \xo{SOLVE_LINEAR} & $A,b\to x$, $Ax=b$ & Square $A$, nonzero determinant, compatible vector and division-capable scalar sort.\\
96 \xo{FIN_SUM} & Int $a,b$, captures $\to S$ & Inclusive bounds. Payload \xo{Type(S) Scope(body)}; index binder last. Empty interval gives additive identity.\\
97 \xo{FIN_PRODUCT} & Int $a,b$, captures $\to S$ & Same payload/binding; empty interval gives multiplicative identity. Terms dimensionless in this base rule.\\
98 \xo{DIFF} & Point, captures $\to$ RealTerm & Payload \xo{Type(variable) Scope(body)}; variable binder last. Derivative must exist at the point.\\
99 \xo{INTEGRAL} & $a,b$, captures $\to$ RealTerm & Same payload/binder; oriented finite Riemann integral. Integrability is required.\\
100 \xo{LIMIT} & Point, captures $\to$ RealTerm & Payload \xo{Type(variable) O(side) Scope(body)}. Side 0 two-sided, 1 from below, 2 from above. Unique finite limit required.\\
101 \xo{ODE} & $t_0,s_0,t$, captures $\to$ state & Payload \xo{Q(interval_left) Q(interval_right) Scope(F)}. Closed interval includes $t_0,t$; time binder then state last. Unique continuous integral-equation solution required.\\
102 \xo{SOLVE_ROOT} & Captures $\to$ RealTerm & Payload \xo{Q(left) Q(right) Scope(f)}; binder $x$ last. Exactly one real root in the open rational interval; existence and uniqueness required.\\
112 \xo{ASA_NA} & $x,V,A,N,B\to(a,n,h,y)$ & Five Bits$(w)$ operands; result has word components $a,n,y$ and Int $h$. Whole-word absorption below.\\
113 \xo{JK} & $q,J,K,V\to$ Bits$(w)$ & Four same-width words, reading the same old $q$; complement restricted to width.\\
\end{longtable}}

\textbf{Calculus binders and solution domain.} Captures bind in operand order,
followed by the displayed integration/differentiation variables. Each body has
one export. Body and variable dimensions determine derivative/integral units.
\xo{ODE} supports a scalar, uniform vector, framed vector, or finite nested tuple
of those dimensioned components; every scalar result leaf has sort
\xo{RealTerm}. Initial algebraic values are coerced componentwise while
preserving all frame annotations. $F$ has the same structure and identical
nominal frame on each framed vector, with each scalar leaf dimension equal to
its state dimension minus the time dimension. Rotating-frame derivatives still
require the appropriate terms in $F$; the type supplies none. Interval rationals
use the same declared time unit as $t_0$ and $t$.

Specifically, \xo{ODE} denotes the unique continuous $s$ on the declared closed
interval $I$ such that, componentwise,
\[
 s(u)=s_0+\int_{t_0}^{u}F(v,s(v))\,dv\qquad(u\in I),
\]
with a Riemann-integrable integrand along $s$. It is differentiable where the
composite right-hand side is continuous; a derivative at a coefficient jump
is not required. The derived name \xo{FLOW} uses this same declaration.
State-dependent discontinuities still require existence/uniqueness under the
event policy. A local theorem does not certify the whole interval. General root
and ODE declarations are not guaranteed solution algorithms. Unbounded sums,
improper integrals, complex roots and distributional derivatives are outside
these opcodes.

\textbf{Dimensions.} If operands have dimension vectors $d_A,d_B$, then
\xo{MUL}, \xo{DOT}, \xo{CROSS} and \xo{MATMUL} return $d_A+d_B$;
\xo{DIV} returns $d_A-d_B$. \xo{NORM} retains dimension. For a uniformly
dimensioned $n\times n$ matrix, determinant has dimension $nd_A$, inverse
$-d_A$, and \xo{SOLVE_LINEAR} has dimension $d_b-d_A$. Addition, subtraction
and order require equal dimensions. Heterogeneous Jacobians use dimensioned
tuples/blocks and explicit contraction templates, not one falsely uniform
matrix type. Dimensioned powers/roots must establish the annotated output
dimensions from their degree at admission; a variable degree cannot contradict
a fixed annotation. Dimensionless inputs avoid that degree/unit issue.

The retained kernel operators have literal equations
\[
 a=x\land A\land V,\qquad n=a\land N,\qquad h=\pop(n\land B),
\]
\[
 y=\begin{cases}0,&h>0,\\n,&h=0,\end{cases}
 \qquad\operatorname{ASA\_NA}(x,V,A,N,B)=(a,n,h,y),
\]
\[
 \operatorname{JK}(q,J,K,V)=
 ((J\land\neg_w q)\lor(\neg_w K\land q))\land V.
\]
Here $\neg_wq=(2^w-1)\oplus q$. The second mask set receives the first stage's
$y$; both traces remain exported when required. $J,K$ use the same old state
and final $y$. Absorption zeroes the whole stage output and does not silently
substitute a hold rule. The state table supplies the delay.

\subsection{Named templates and the proof-record interface}
Named templates require explicit typed parameters and bodies. Labels alone
provide no operator implementation. \xo{WRAP} is
$x-P\lfloor x/P\rfloor$ for $P>0$, paired with integer winding;
\xo{PHI} is $(1+\sqrt5)/2$; \xo{FIBONACCI} has initial $(0,1)$ and delayed
pair update $(a,b)\mapsto(b,a+b)$. \xo{NORMALIZE} requires nonzero norm.
\xo{GRADIENT}/\xo{JACOBIAN} assemble coordinate \xo{DIFF} terms.
Polynomial/Horner and Chebyshev templates preserve their finite coefficients
and degrees, with $T_0=1,T_1=x,T_{n+1}=2xT_n-T_{n-1}$.
Legendre, gravity, frame, third-body, pressure/shadow and station templates retain
their source equations and singular/branch domains. \xo{RK4} is the finite rule
in Section~\ref{sec:exact-core}, not the continuous \xo{ODE} declaration.
\xo{OTAN2}, both UGTS key layouts and the source geometry remain separate
templates with their literal source identities. \xo{OTAN2} is not \xo{ATAN2}.
Modular keys do not replace absolute coordinates, time or sample bytes.

The strategy and proof records have this exact outer order:
\begin{lstlisting}
Strategy = U(version) Name(proof_system) Seq_Rule(rules)
           B(priority_policy) B(cost_model) B(resource_policy)
           Seq_Proposal(history)
Rule = Name(label) B(typed_statement) B(proof_object)
Proposal = U(parent_version) U(new_version)
           B(old_canonical_graph) B(new_canonical_graph)
           B(statement) B(proof_object) B(cost_argument)
\end{lstlisting}
The named proof-system/policy profile resolves the \xo{B} contents; they are
not executable arbitrary source. The baseline with empty rules, policy and
history admits identity-only execution strategy. XOP-R1 defines the interface,
not a completed proof language or checker. Uninterpreted policy bytes confer
no authority to admit a nonidentity rewrite.

For $\Gamma\vdash e\simeq e':T$, domains under $\Gamma$ and denotations on
those domains must be identical. The proof interface has these rule families:
{\small
\begin{longtable}{@{}P{.31\linewidth}P{.64\linewidth}@{}}
\toprule Rule family & Required premise\\\midrule\endfirsthead
\toprule Rule family & Required premise (continued)\\\midrule\endhead
\bottomrule\endfoot
\xo{ASSUMPTION} & The exact proposition occurs in the typed environment.\\
\xo{REFL}, \xo{SYM}, \xo{TRANS} & Typed equality and compatible domain judgements.\\
\xo{CONGRUENCE} & Equal corresponding operands, same operator/bindings and preserved definedness.\\
\xo{SUBSTITUTION}, \xo{BETA} & Capture-avoiding typed substitution; function/body domain preserved.\\
\xo{EXPAND} & Exact catalog/template definition, including its branches and domains.\\
\xo{INTEGER_RING}, \xo{RATIONAL_FIELD} & Exact laws with every nonzero-denominator obligation retained.\\
\xo{ORDER}, \xo{ROOT_BRANCH} & Required sign, range and unique-root facts.\\
\xo{ANALYTIC_LEMMA} & Explicit derivative/integral/limit hypotheses and proof in the accepted calculus.\\
\xo{STATE_SIMULATION} & Initial relation, step-domain equivalence and equal declared observations for every admitted input.\\
\end{longtable}}
A theorem identifier must resolve to an accepted theorem with discharged
substitutions and premises. A named or absent checker cannot admit nonidentity
proofs. Strategy versions are immutable; the previously accepted checker reviews
a proposal. Replacing that checker creates a new trust boundary and cannot be
certified merely by the proposal's own changed rules. Hashes identify bytes,
not soundness or termination.

Memoization binds complete operator/dependency identity, arguments, sample
snapshot and state step. Resource outcomes and traces also count when promised
by the interface. A decreasing natural-valued measure proves termination of
that rewrite phase; a finite search budget proves only bounded search. Neither
establishes global optimality, speed, bounded limb growth or physical accuracy.
The causal and conditional replacement proofs remain those of
Section~\ref{sec:exact-core}; no new algorithm validation is asserted here.
\endgroup
